\section*{Chapitre \ref{ch:b2:wilson-cycle} --- Massifs anciens et cycle de Wilson}

\begin{solution}{exo:b2:wilson-cycle:1}
(a) $600 - 400 = \qty{200}{Ma}$, depuis le rifting des marges jusqu'à la dernière collision.
(b) $480 - 400 = \qty{80}{Ma}$. (c) L'orogenèse calédonienne est un ensemble de phases étalées sur
près de 200~Ma, et chaque phase a commencé à des moments différents selon les lieux~: le mot
désigne un océan et sa chaîne, non un moment. Un âge se donne par l'échelle stratigraphique ou en
Ma.
\end{solution}

\begin{solution}{exo:b2:wilson-cycle:2}
(c) le rift est-africain, rifting~; (a) la mer Rouge, un océan jeune~; (d) l'Atlantique, un océan large
à marges passives~; (e) la marge pacifique de l'Amérique du Sud, subduction et
fermeture d'un océan~; (b) l'Himalaya, collision~; (f) la Bretagne, une chaîne aplanie, après quoi
un rifting peut reprendre.
\end{solution}

\begin{solution}{exo:b2:wilson-cycle:3}
(a) Après le dépôt des couches du Silurien inférieur, donc après \qty{443.1}{Ma}, et avant le
dépôt des couches du Dévonien supérieur, donc avant \qty{358.9}{Ma}, la fin du Dévonien.
(b) Moins de $443.1 - 358.9 = \qty{84.2}{Ma}$~; en fait moins encore, puisque les couches inférieures sont plus
jeunes que la base du Silurien et les couches supérieures plus anciennes que son sommet. (c) Dépôt des
turbidites grises dans l'océan~; plissement et clivage lors des collisions calédoniennes~; soulèvement et
érosion jusqu'à une surface continentale irrégulière~; dépôt de brèches et de conglomérats rouges~; légère
bascule lors de l'orogenèse varisque~; l'érosion actuelle, qui entaille la falaise.
\end{solution}

\begin{solution}{exo:b2:wilson-cycle:4}
(a) Les deux~: un bloc de croûte continentale séparé des continents par des océans pendant
qu'il traversait l'Iapetus, et, une fois accolé, un bloc délimité par des failles dont l'histoire diffère de celle de ses
voisins. (b) Aucun des deux~: un continent, et non un bloc ajouté à un continent. (c) Un \omterm{def:b2:wilson-cycle:terrane}{terrane}, et non un
\omterm{def:b2:wilson-cycle:terrane}{microcontinent}~: sa croûte est celle d'un arc, non un fragment de continent. (d) Aucun des deux~: c'est
la limite entre \omterm{def:b2:wilson-cycle:terrane}{terranes}, la trace de l'océan refermé.
\end{solution}

\begin{solution}{exo:b2:wilson-cycle:5}
(a) $3000 - 2600 = 400$, $2600 - 1500 = 1100$, $1500 - 875 = 625$, $875 - 600 = 275$, $600 -
262.5 = 337.5$~; moyenne $2737.5/5 = \qty{547.5}{Ma}$. (b) $1100/275 = 4.0$. (c) $(1500 -
262.5)/3 = \qty{412.5}{Ma}$. La moyenne varie d'un quart lorsque les points douteux sont retirés, et les
écarts varient d'un facteur quatre~: une récurrence grossière, non une période.
\end{solution}

\begin{solution}{exo:b2:wilson-cycle:6}
(a) $3000/(70 - 25) = \qty{66.7}{Ma}$. (b) $3000/70 = \qty{42.9}{Ma}$. (c) Une fois la dorsale subduite, le
plancher océanique qui reste appartient à la plaque éloignée, qui porte le continent lointain~: la vitesse de
fermeture est la vitesse de cette plaque par rapport à la fosse.
\end{solution}

\begin{solution}{exo:b2:wilson-cycle:7}
(a) $700/8 = \qty{87.5}{km}$~; $700/30 = \qty{23.3}{km}$. (b) La chaîne varisque~: à
\qty{30}{\celsius} par kilomètre, la croûte moyenne fond~; à \qty{8}{\celsius} par kilomètre
la température de fusion se situe sous la base d'une croûte même épaissie. (c) Le gradient faible décrit des
roches entraînées vite vers le bas dans une zone de subduction, avant de pouvoir se réchauffer. Les roches qui restent
profondément enfouies dans une croûte épaissie se réchauffent vers un géotherme normal en quelques Ma, et peuvent
alors atteindre la fusion à plus faible profondeur.
\end{solution}

\begin{solution}{exo:b2:wilson-cycle:8}
(a) $\tau\ln(3/0.3) = 26.3\times2.303 = \qty{60.6}{Ma}$. (b) $290 - 60.6 = \qty{229}{Ma}$, au
Trias supérieur. (c) $\tau = (290 - 251.9)/\ln10 = \qty{16.5}{Ma}$. (d) L'effondrement extensif de la
croûte épaissie, qui l'amincit et abaisse directement sa surface~; également une
vitesse d'érosion plus forte, un climat plus humide ou des roches plus tendres, ou une altitude de départ plus basse.
\end{solution}

\begin{solution}{exo:b2:wilson-cycle:9}
(a) La suture, trace de l'Iapetus, se trouve à terre en Grande-Bretagne et en Irlande, à l'est
de l'Atlantique~; la croûte de Laurentia, venue du côté lointain de l'Iapetus, se trouve près de Rockall, du
côté européen. Le nouveau rift s'est ouvert à l'ouest de l'ancienne suture~: il s'est ouvert près d'elle, non le long d'elle.
(b) Au nord de la suture, socle et couverture de la marge laurentienne~; au sud, ceux d'Avalonia,
aux histoires différentes et, au Paléozoïque inférieur, aux faunes différentes. (c) Le rifting
est survenu après 80~Ma, soit plus que les 30 à 50~Ma au bout desquels une zone refroidie peut être devenue
résistante~; pourtant l'Atlantique central a bien suivi les Appalaches. La chaleur d'une chaîne s'estompe~; ses
failles, fabriques et compositions demeurent et peuvent encore guider un rift.
\end{solution}

\begin{solution}{exo:b2:wilson-cycle:10}
(a) $70 - 40/2 = \qty{50}{mm}$ par an. (b) $\qty{1500}{km}/\qty{50}{km.Ma^{-1}} =
\qty{30}{Ma}$. (c) L'océan se rétrécit à $70 - 40 = \qty{30}{mm}$ par an~: $3000 - 30\times30 =
\qty{2100}{km}$. (d) Une lithosphère jeune et chaude arrive alors à la fosse, et aucune plaque vieille et froide
ne se déshydrate plus sous l'arc~: le magmatisme de l'arc change ou s'arrête.
\end{solution}

\begin{solution}{exo:b2:wilson-cycle:11}
(a) $h = \Delta T(\rho_m - \rho_c)/\rho_m = 25\times600/3300 = \qty{4.5}{km}$. (b) Il faut enlever toute la
surépaisseur, \qty{25}{km}, soit 5.5 fois la hauteur de la chaîne. (c) $2700\times9.8\times25\,000 =
\qty{6.6e8}{Pa}$, soit environ \qty{0.66}{GPa}. (d) À environ \qty{0.7}{GPa} et aux températures de
la racine d'une croûte épaissie, plusieurs centaines de degrés, le faciès à amphibolites ou celui à
granulites~: les gneiss et les \omterm{def:b2:crustal-melting:migmatite}{migmatites} des vieux massifs.
\end{solution}

\begin{solution}{exo:b2:wilson-cycle:12}
(a) $S = 30\times10 = \qty{300}{km}$~; $W_0 = 300 + 300 = \qty{600}{km}$~; $T = 35\times600/300 =
\qty{70}{km}$. (b) $h = 35\times600/3300 = \qty{6.36}{km}$~; $r = 4.5h = \qty{28.6}{km}$.
(c) $\Delta F = 2700\times9.8\times6364\times(35\,000 + 3182 + 14\,318) =
\qty{8.8e12}{N.m^{-1}}$, 1.4 fois celle du Tibet. (d) Une telle chaîne est plus lourde que le Tibet, qui
s'étale déjà~: \omterm{def:b2:himalaya-andes:collapse}{effondrement gravitaire}, par extension de ses régions élevées et chevauchements sur ses bords.
\end{solution}

\begin{solution}{pb:b2:wilson-cycle:1}
\textbf{1.} $720/(540 - 500) = \qty{18}{km}$ par Ma, soit \qty{18}{mm} par an.
\textbf{2.} \qty{36}{mm} par an~: une dorsale à expansion lente (de 20 à \qty{55}{mm} par an).
\textbf{3.} $36\times(580 - 540) = \qty{1440}{km}$.
\textbf{4.} $36\times(580 - 500) = \qty{2880}{km}$.
\textbf{5.} Expansion symétrique à vitesse constante, à partir d'une largeur nulle à \qty{580}{Ma}, et aucune
subduction avant \qty{500}{Ma}.
\textbf{6.} \qty{80}{Ma}, aux deux bords de l'océan, contre les marges continentales. C'est la
lithosphère océanique la plus froide, la plus épaisse et la plus dense de l'océan, la plus prompte à plonger.
\textbf{7.} $64 - 36 = \qty{28}{mm}$ par an.
\textbf{8.} $2880/28 = \qty{102.9}{Ma}$ après \qty{500}{Ma}~: à \qty{397.1}{Ma}.
\textbf{9.} La dorsale, à \qty{1440}{km} de la fosse, s'en approche à $64 - 18 =
\qty{46}{mm}$ par an~: après $1440/46 = \qty{31.3}{Ma}$, à \qty{468.7}{Ma}.
\textbf{10.} $2880 - 28\times31.3 = \qty{2004}{km}$.
\textbf{11.} $2004/40 = \qty{50.1}{Ma}$ plus tard~: à $468.7 - 50.1 = \qty{418.6}{Ma}$.
\textbf{12.} Dans l'intervalle de 425 à \qty{400}{Ma} de la collision scandienne.
\textbf{13.} $40\times8 = \qty{320}{km}$.
\textbf{14.} $W_0 = 400 + 320 = \qty{720}{km}$~; $T = 35\times720/400 = \qty{63.0}{km}$.
\textbf{15.} $h = 28\times600/3300 = \qty{5.09}{km}$~; $r = 4.5\times5.09 = \qty{22.9}{km}$.
\textbf{16.} $\Delta F = 2700\times9.8\times5091\times(35\,000 + 2545 + 11\,455) =
\qty{6.60e12}{N.m^{-1}}$, environ \qty{8}{\percent} de plus que celle du Tibet.
\textbf{17.} Pour cesser de s'élever et de s'étaler~: extension et failles normales dans son intérieur élevé,
chevauchements sur ses bords, comme le fait le Tibet aujourd'hui.
\textbf{18.} $f = 600/3300 = 0.182$~; $k = 2.303/11 = 0.209$ par Ma~; $\tau = 1/(kf)
= \qty{26.3}{Ma}$.
\textbf{19.} $3.0\,e^{-30/26.3} = \qty{0.96}{km}$.
\textbf{20.} $26.3\ln(3.0/0.5) = \qty{47.1}{Ma}$ après \qty{410}{Ma}~: à \qty{362.9}{Ma}.
\textbf{21.} La croûte en excès vaut $h/f$~: $3.0/0.182 = \qty{16.5}{km}$ se réduit à $0.5/0.182 =
\qty{2.75}{km}$, donc on enlève \qty{13.75}{km} de roche.
\textbf{22.} $2700\times9.8\times13\,750 = \qty{3.6e8}{Pa}$, soit environ \qty{0.36}{GPa}.
\textbf{23.} Oui~: au Dévonien supérieur, la chaîne du modèle est tombée à environ \qty{1}{km} d'altitude
moyenne et continue de baisser, avec du relief entre ses vallées et ses crêtes~; des graviers rouges peuvent
s'accumuler dans ses creux bien avant \qty{358.9}{Ma}, comme à Siccar Point, où ils reposent sur une surface
dotée de relief.
\textbf{24.} \textbf{\qty[text-series-to-math=true]{2880}{km}}, à \qty{500}{Ma}, au moment où
la subduction a commencé.
\end{solution}
